% ARCHIVED at v3.97 (2026-09-24) -- author: "reclaim why we do the retest the K1/2 in UAV-pillars: they are the
% results of the Pareto optimal from d3il-avoiding, so we run it again on the UAV and test what else will happen".
% Verbatim, not part of the build.

% ----
not retrained (\autoref{sec:setup:tasks:uav}). What is read here is therefore not a new benchmark but a
check of the vehicle and its tracking controller: the configurations that win on D3IL-avoiding --- the two
average-velocity models at one and two network evaluations, the operating points of
\autoref{tab:avoiding-conclusion} --- are flown, before and after per-step projection, and their flown
outcome is set beside their table outcome at the same protocol, five training seeds $\times$ two episodes
per geometry, averaged over the three geometries.

% ----
UAV-pillars reaffirms D3IL-avoiding on the quadrotor. Its demonstrations, planner, projection and constraint
rows are those of D3IL-avoiding, and only the plant that executes each commanded position is replaced
(\autoref{sec:method:env:uav}).

% ----
What the quadrotor and its controller change against the manipulator is twofold.

% ----
\guard{The diffusion baseline and instantaneous-velocity matching were not flown on this scene; the dominance
over the baseline is a result of D3IL-avoiding.}

% ----
  \item[UAV-pillars: transfer of the D3IL-avoiding result to the quadrotor.] The declared constraints transfer exactly
    and the two average-velocity models keep their order.
